\section{The signed old-stress response to young affine forcing}

Fix a real compact interval $K\Subset(0,1/4)$ and $c=\nu^{-2}$.
Cell $L^2$ norms use the unnormalized measure $dq\,d\eta$ on
$(-c/2,c/2)\times\mathbb R$; stress averages contain $c^{-1}$.
Velocity components are ordered $(r,\phi,z)$, with zero azimuthal velocity.
The vorticity sources are centrally even and have zero circulation.
Let
\[
 V_G(x,s)=a(s)^{-2}e^{-|x|^2/(4a(s))}Jx,\qquad a(s)=1+\nu s,
 \qquad V_j=V_G(q-jc,\eta,s).
\]
Write $E_s=\operatorname{curl}V_G(\cdot,s)$, and let $H_{s,c}$ and
$R_{s,c}$ be the periodizations of $rE_s'(r)e^{2i\vartheta}$ and
$rE_s'(r)+2E_s(r)$ respectively. The Gaussian defines the source;
the evolution operator $L_{\nu,c}(s)$ uses the actual cylinder train.
Take $\chi_f(s)=\widehat\chi_f(\nu s)$ with $\widehat\chi_f=1$ on
$[0,1]$ and support in $[0,1.1)$, and take $\chi_o=0$ for
$s\le1.2/\nu$ and $\chi_o=1$ for $s\ge1.5/\nu$. Set
\[
\begin{gathered}
 I_T(M)=\int_0^\infty\chi_o(s)s^{-3/2}
       \frac1c\int_{\rm cell}(2M_{zz}-M_{rr})\,dq\,d\eta\,ds,
 \qquad \ell_\mu(M)=\frac{\epsilon^3I_T(M)}{8(\mu+3/2)},\\
 N_\nu=e^{\lambda/\nu}\nu^{-\kappa},\quad
 \gamma=\operatorname{Re}\lambda,\quad \eta_\nu=\nu+\delta_c/\nu,
 \quad \mathcal A_\nu=|N_\nu|\nu^{3+2\mu}/c,\quad
 A_*=\epsilon^3\mathcal A_\nu.
\end{gathered}
\]

Fix $0<a<\gamma/16$. Let $\mathscr B_a$ consist of measurable real
trace-free matrices $B(s)$, with no mixed azimuthal entries, supported
in $0<s<1.1/\nu$, with norm
\[
 \|B\|_a=\mathop{\rm ess\,sup}_{s>0}e^{-as}s^{-3/2}|B(s)|.
\]
For $x_j=(q-jc,\eta)$ define the full affine force
\[
 f_B=-P_0\sum_j[(B_{\rm mer}x_j)\cdot\nabla V_j+B V_j].
\]
Decompose
\[
 B_{\rm mer}=\beta I+\begin{pmatrix}d&e\\e&-d\end{pmatrix}
       +\omega\begin{pmatrix}0&-1\\1&0\end{pmatrix},
 \qquad B_{\phi\phi}=-2\beta.
\]
The antisymmetric contribution is a periodic gradient, annihilated
by $P_0$. The vorticity source is
\[
 g_B=-\tfrac12(d-ie)H_{s,c}-\tfrac12(d+ie)\overline{H_{s,c}}
          -\beta R_{s,c}.
\]
Let
\[
 \xi_B(s)=s^{-\mu}\int_0^s t^\mu\mathcal U_{\nu,c}(s,t)g_B(t)\,dt,
 \qquad b_B=K_c\xi_B,\qquad M_B=v\otimes b_B+b_B\otimes v.
\]
Since $g_B=O(s^{3/2})$ near zero, $\xi_B=O_{X_{m,c}}(s^{5/2})$;
this causal column has finite-energy entrance. Define its limiting amplitude
\[
 a_0(B)=-\tfrac12\ell_+(H_0)
              \int_0^\infty t^\mu e^{-\lambda t}(d(t)-ie(t))\,dt.
\]

\begin{proposition}[Uniform affine-source return]
There is a canonical modulus $\omega_{\rm can}(\nu)\to0$ such that,
uniformly on $K$ and $\mathscr B_a$,
\begin{equation}\label{flat-return:class-amplitude}
 |a_{\nu,+}(B)-a_0(B)|\le C_{K,a,\Sigma}\eta_\nu\|B\|_a,
\end{equation}
where $a_{\nu,+}$ is the frozen-source amplitude of the young-source theorem.
For every fixed $0<\gamma'<\gamma$,
\begin{align*}
 \biggl|\ell_\mu(M_B)
  +\frac{3\pi}{16}\frac{\epsilon^3\nu^{3+2\mu}}c
       \operatorname{Re}[i a_0(B)N_\nu J(\mu)]\biggr|
 \le C_{K,a,\Sigma}A_*\bigl[&\omega_{\rm can}+\eta_\nu
       +P_\nu e^{-\gamma'/\sqrt\nu}\bigr]\|B\|_a.
\end{align*}
For $k=0,1$ one also has
\begin{equation}\label{flat-return:all-moments}
 \int_0^\infty\chi_o(s)s^{-3/2-k/2}\frac1c
       \int_{\rm cell}|M_B(s,q,\eta)|\,dq\,d\eta\,ds
 \le C_{K,a,\Sigma,k}\mathcal A_\nu\nu^k\|B\|_a.
\end{equation}
\end{proposition}
\begin{proof}
The explicit amplitude is
\[
 a_{\nu,+}(B)=\int_0^{1.1/\nu}t^\mu F_+(t,0)^{-1}
                            \ell_{+,\nu t,c}g_B(t)\,dt,
 \qquad F_+(t,0)^{-1}=e^{-\lambda t/(1+\nu t)}(1+\nu t)^{-\kappa}.
\]
The direct/dual Gaussian comparisons and heat scaling give
\[
 \ell_{+,\nu t,c}H_{t,c}=\ell_+(H_0)+O(\nu t+c^{-2}),\qquad
 \ell_{+,\nu t,c}\overline{H_{t,c}},\ \ell_{+,\nu t,c}R_{t,c}=O(c^{-2}),
\]
\[
 |F_+(t,0)^{-1}-e^{-\lambda t}|
       \le C\nu(t+t^2)e^{-\gamma t/3},\qquad 0\le t\le1.1/\nu.
\]
The resulting error has integrable majorant
$C(\eta_\nu+\nu(t+t^2))t^{\mu+3/2}e^{-(\gamma/3-a)t}\|B\|_a$,
proving \eqref{flat-return:class-amplitude}. The radial and opposite
harmonic sources have zero limiting projection; their source norms
need not be small.

Apply the young-source theorem with $\alpha=1.1$. Its tested
source norm is at most $C\|B\|_a$, since
$t/(1+\nu t)\ge t/2.1$ on its support. Thus
\eqref{eq:young-later-output} transports the finite-band error to the
whole old output without a polynomial loss. The real form of
\eqref{eq:young-old-anisotropy} gives
\[
 I_T(M_B; s\ge u_o)=\frac{3\pi}{2c}(\mu+3/2)\nu^{3+2\mu}
      \operatorname{Im}[a_{\nu,+}(B)N_\nu J(\mu)]
      +O_K(\mathcal A_\nu(\eta_\nu+\omega_{\rm can}))\|B\|_a,
 \quad u_o=\nu^{-3/2}.
\]
The earlier-old part has bound
$\mathcal A_\nu P_\nu e^{-\gamma'/\sqrt\nu}\|B\|_a$ by the full
forward gain estimate. The old cutoff acts on the full tensor.
The source ends before that cutoff begins, so the old moment law is
homogeneous; both weighted endpoints were retained in the canonical
calculation. Substitute \eqref{flat-return:class-amplitude} and divide
by $8(\mu+3/2)/\epsilon^3$ to obtain the scalar estimate.

For the absolute moments use
$s^{-k/2}=\nu^k\theta^{k/2}(1-\nu\theta)^{-k/2}$.
On $\theta\le b$, apply \eqref{eq:ct-absolute-impulse} with
$p=\mu+3/2+k/2$. On $b\le\theta\le\theta(u_o)$ the denominator
is bounded and the exponential long-clock factor absorbs
$\theta^{k/2}$. The earlier-old exponential bound absorbs its fixed
polynomial age factors. Apply these estimates to both bare columns
and to the uniform frozen-pair error, using the state-at-$H$ bound
\eqref{eq:young-later-output}. This proves the absolute estimates on
the whole source unit ball.
\end{proof}

\begin{corollary}[Unit signed coefficient]\label{flat-return:coefficient}
Put
\[
\begin{gathered}
 D_0=\operatorname{diag}(12,-3,-9),\qquad
 A_\mu(s)=\frac{\chi_f(s)s^{5/2}}{\mu+5/2}D_0,\qquad
 B_\mu=A_\mu'+\frac\mu s A_\mu=r_\mu(s)D_0,\\
 r_\mu(s)=\chi_f(s)s^{3/2}
      +\frac{\chi_f'(s)s^{5/2}}{\mu+5/2},\qquad
 L(\mu)=\ell_+(H_0)\frac{\Gamma(\mu+5/2)}{\lambda^{\mu+5/2}},
 \qquad \mathcal C(\mu)=\frac{63\pi i}{128}L(\mu)J(\mu).
\end{gathered}
\]
Let $(f_\mu,\xi_\mu,b_\mu,M_\mu)=(f_{B_\mu},\xi_{B_\mu},b_{B_\mu},M_{B_\mu})$.
For some $b>0$,
\[
 \left|\ell_\mu(M_\mu)-2\epsilon^3\nu^{5+2\mu}
                 \operatorname{Re}[\mathcal C(\mu)N_\nu]\right|
 \le C_KA_*[\omega_{\rm can}+\eta_\nu+
               P_\nu e^{-b/\nu}+P_\nu e^{-\gamma'/\sqrt\nu}].
\]
\end{corollary}
\begin{proof}
Here $\beta=3r_\mu/2$, $d=21r_\mu/2$ and $e=\omega=0$, so
\begin{equation}\label{flat-return:source}
 g_\mu=-\frac{21}{4}r_\mu(H_{s,c}+\overline{H_{s,c}})
              -\frac32r_\mu R_{s,c}.
\end{equation}
This retains the trace of the $(r,z)$ block $3r_\mu$ and hoop entry $-3r_\mu$.
The unit $\mathscr B_a$ norm is bounded. Its cutoff derivative is
supported above $1/\nu$, so the explicit Laplace integral gives
$a_0(B_\mu)=-(21/4)L(\mu)+O_K(P_\nu e^{-b/\nu})$.
Substitute this into the general scalar formula and use $c^{-1}=\nu^2$.
\end{proof}
